\chapter{開發過程中所使用的 AI 工具與對話紀錄}

在本專案的開發過程中，我們使用了對話式人工智慧工具（Google Gemini）作為輔助開發手段，
以協助解決實作過程中所遇到的技術問題。其應用範圍涵蓋資料結構設計、語法解析邏輯、
有限狀態機（FSM）建構，以及程式碼優化等面向。

AI 工具主要用於提供設計建議、比較不同實作方案，以及協助釐清複雜問題的思路，
最終決策與實作仍由開發者自行完成與驗證，以確保系統行為之正確性。

以下列出開發過程中具有代表性的對話紀錄摘要（完整內容請參見 \lstinline{AI_prompt} 資料夾中的 PDF 檔案）：

\begin{enumerate}
    \item SIC Opcode 查表資料結構比較
          \href{./AI_prompt/SIC-Opcode-查表資料結構比較.pdf}{SIC-Opcode-查表資料結構比較.pdf}

    \item SIC/XE Format 2 指令暫存器限制
          \href{./AI_prompt/SIC-XE-Format-2-指令暫存器限制.pdf}{SIC-XE Format 2 指令暫存器限制.pdf}

    \item SIC/XE 指令 C 語言解析器設計
          \href{./AI_prompt/SIC-XE-指令-C-語言解析器.pdf}{SIC-XE-指令-C-語言解析器.pdf}

    \item 狀態機設計與實作優化
          \href{./AI_prompt/狀態機-枚舉與實作優化.pdf}{狀態機-枚舉與實作優化.pdf}
\end{enumerate}
